% Current FORM implementation; shared by both standalone references.
% New displays are unnumbered to preserve the earlier equation numbering.
\sessionheading{Degree-bounded FORM exponential: current implementation}
\label{sec:degree-bounded}
The 16 September implementation in \code{_build_form_program} evaluates the
same finite exponential as the auxiliary-variable recurrence above, but avoids
products that cannot survive the requested cutoff. Let $N$ denote that cutoff
(called $D$ in the earlier formulas), and write
\[
 A=\texttt{itotal}=\sum_{d=2}^{N}A_d t^d,\qquad
 [P]_N=\sum_{d\leq N}[t^d]P\,t^d,\qquad M=\lfloor N/2\rfloor.
\]
The coefficients $A_d$ contain Laurent monomials in $y,u$ and formal gauge
characters. They contain no $t$. All powers of $t$ are nonnegative and the
exponent starts at degree two. The target before gauge projection is
$[\exp(A)]_N$. The derivation below follows the project's generated program.

\textbf{Step 1: construct the derivative series and plethystic exponent.}
The expression \code{J} is the truncated product of geometric series for
$(1-t^3y)^{-1}(1-t^3/y)^{-1}$. Substituting the vector and hyper letters into
\code{itotal} gives
\[
 A=\left[\sum_{j=1}^{M}\frac{1}{j}
 \left(f_V(t^j,y^j,u^j)\sum_{p\in V}C_p(j)
 +f_H(t^j,y^j,u^j)\sum_m N_m\prod_{p\in m}C_p(j)\right)\right]_N.
\]
Here $j$ is the Adams index. It differs from the outer Taylor-power index $i$
below. Integer multiplicities $N_m$ are unchanged by Adams substitution.
The bound on $j$ follows from $2j\leq N$; the derivative rectangle uses
\code{order // 3}. The declaration \code{t(:N)} remains a final degree guard.

\textbf{Step 2: organize the exponent and initialize the active power.}
\begin{lstlisting}
Bracket t;
.sort
Skip J,itotal;
L I=z*itotal;
.sort
\end{lstlisting}
\code{Bracket t} groups equal $t$ powers. A reference \code{itotal[t^d]}
returns $A_d$, so multiplication must restore \code{t^d} explicitly.
\code{Skip J,itotal;} preserves those expressions, including the bracket
information, for the current module. They remain readable on the right-hand
side. Repeat Skip after each module boundary; it is not a permanent freeze
\cite[entries ``bracket'', ``skip'', and Chapter 9]{form-current}.

Initially $I=zA$. The symbol $z$ marks the one Taylor order still eligible
to generate the next order; unmarked terms have already been accumulated.
Neither $z$ nor the temporary symbol $w$ is a physical fugacity.

\clearpage
\sessionheading{One exponential step, line by line}
For example, at cutoff $N=8$ the following is the actual emitted loop and
finalization block; Python substitutes the integer bounds before running FORM:
\begin{lstlisting}
#do i=2,4
  Skip J,itotal;
  #do k=2,6
    if (count(t,1) == `k');
      id z=1+w*sum_(idx1,2,8-`k',itotal[t^idx1]*t^idx1)/`i';
    endif;
  #enddo
  id z=1;
  id w=z;
  .sort:step `i';
#enddo

L result=1+I;
id z=1;
.sort
Print result;
.end
\end{lstlisting}
\textbf{Step 3: select the current term's degree.}
The inner preprocessor loop creates one branch for each $k=2,\ldots,N-2$.
\code{count(t,1)} tests the current term's $t$ degree. An unmarked term is
unchanged because it has no $z$ to replace. On a marked term $zct^k$, the
substitution has the exact effect
\[
 zct^k\ \longmapsto\ ct^k+
 \frac{w c}{i}\sum_{d=2}^{N-k}A_d t^{k+d}.
\]
The first branch retains the old Taylor contribution. The second produces
the next contribution, dividing by $i$ to turn $(i-1)!$ into $i!$.
Its summation bound is the condition $k+d\leq N$, enforced before forming
the product rather than only when an excessive power is normalized away.

\textbf{Step 4: protect newly generated terms.}
New terms carry $w$, not $z$. A later $k$ branch in the same outer iteration
can encounter one of their new degrees, but its \code{id z=...} cannot expand
them again. Reusing $z$ immediately would allow a single Taylor step to
multiply a newly produced term more than once.

\textbf{Step 5: finalize old terms, then mark the new ones.}
\code{id z=1;} removes any surviving old marker, including terms at degrees
$N-1$ and $N$. Then \code{id w=z;} makes only the new generation active.
Reversing this order would erase the new marker. Finally
\code{.sort:step `i';} collects equal terms and ends the module; the step
label is diagnostic. The next iteration reinstates Skip for its own module.

\clearpage
\sessionheading{Why the bounds and factorials are exact}
For $2\leq i\leq M$, induction on the previous step gives
\[
 \begin{aligned}
 I_1&=zA,\\
 I_2&=A+z[A^2/2!]_N,\\
 I_3&=A+[A^2/2!]_N+z[A^3/3!]_N,\\
 I_i&=\sum_{r=1}^{i-1}[A^r/r!]_N+z[A^i/i!]_N.
 \end{aligned}
\]
Only the last summand has a marker. Step $i+1$ retains it without a marker
and multiplies it by $A/(i+1)$, keeping precisely the products of total degree
at most $N$. This proves the invariant with exact rational coefficients.

\textbf{Why the inner loop stops at $N-2$.}
Every term of $A$ has degree at least two, so an old term of degree $k$ can
produce a retained new term only when $k+2\leq N$. Terms of degree $N-1$ and
$N$ are already present and are retained by \code{id z=1;}. Thus $N-2$ is a
bound on terms that can be multiplied again, not a reduced output cutoff.

\textbf{Why the outer loop stops at $\lfloor N/2\rfloor$.}
$A^i$ has minimum degree $2i$. No higher power can contribute. Adding the
vacuum and removing the last marker therefore gives
\[
 \texttt{result}=1+\sum_{r=1}^{M}[A^r/r!]_N=[\exp(A)]_N.
\]
The public API returns $1$ immediately for orders zero and one. Orders two
and three need only $1+A$; Python omits the multiplication loop. These bounds
use the full-index grading; a Hall--Littlewood letter beginning at degree one
would require different bounds.

\textbf{Worked example, $N=6$.} Set $A=a t^2+b t^3$, with commuting exact
coefficients that may themselves be formal characters. Then
\[
 \begin{aligned}
 I_2={}&a t^2+b t^3+
 z\left(\frac{a^2}{2}t^4+ab t^5+\frac{b^2}{2}t^6\right),\\
 I_3={}&a t^2+b t^3+\frac{a^2}{2}t^4+ab t^5+\frac{b^2}{2}t^6
       +z\frac{a^3}{6}t^6.
 \end{aligned}
\]
In step three, only the marked degree-four term can multiply again. The
degree-five and degree-six contributions become unmarked without being lost.
The final coefficient of $t^6$ is $b^2/2+a^3/6$, as in the ordinary exponential.

\textbf{Validation and cache effect.} \code{test/test_form_degree_bound.py}
checks 42 low-order formal expansions (seven cases at six cutoffs) and one
connected bifundamental at order 18 against the former multiplication. The
16 September recorded full run passed 357 tests in 64.117 s, with no skips.
A matched, sequential FORM-only comparison at order 18 took 2.236 s before
and 0.585 s afterward (3.82 times faster), with exactly the same 28,154 formal
terms. This is one measurement, not a universal speedup. Evidence is retained
in \code{output/benchmarks/degree_bounded_form_20260916.json}.
The changed raw program has a new FORM-cache key. Existing rows remain valid
for their original programs, and character-cache entries remain reusable.
